%! Conic Sections — Unit 4, Lesson 4.6 — Group Activity
\documentclass[10pt]{article}
\usepackage{precalculus-article}
\usepackage{precalculus-boxes}
\newcommand{\TallMath}[1]{$\displaystyle #1\rule[-1.4em]{0pt}{3.2em}$}

\begin{document}

\pageheader{Unit 4, Lesson 4.6}{Group Activity}
\namepartnerperiod

\vspace{0.10in}

%----------------------------------------------------------------------
% Tier R
%----------------------------------------------------------------------
\begin{tcolorbox}[colback=white, colframe=black!40,
    title=\textbf{Tier R --- Remediate},
    fonttitle=\bfseries, arc=2mm, left=3mm, right=3mm, top=2mm, bottom=2mm]

\textbf{Directions:} For each equation below, identify the type of conic,
the center or vertex, and the key measurements. No equation-writing is required.

\vspace{6pt}
\textbf{Card A:}\quad $y - 4 = 3(x + 2)^2$

Type: \blank{2.5cm} \quad Vertex: \blank{3cm} \quad $a =$ \blank{1.5cm} \quad Opens: \blank{2cm}

\vspace{10pt}
\textbf{Card B:}\quad $\dfrac{(x-3)^2}{25} + \dfrac{(y+1)^2}{9} = 1$

Type: \blank{2.5cm} \quad Center: \blank{3cm} \quad $a =$ \blank{1.5cm} \quad $b =$ \blank{1.5cm}

\vspace{10pt}
\textbf{Card C:}\quad $(x+4)^2 + (y-2)^2 = 36$

Type: \blank{2.5cm} \quad Center: \blank{3cm} \quad $r =$ \blank{1.5cm}

\vspace{10pt}
\textbf{Card D:}\quad $\dfrac{(x-1)^2}{9} - \dfrac{(y-2)^2}{4} = 1$

Type: \blank{2.5cm} \quad Center: \blank{3cm} \quad $a =$ \blank{1.5cm} \quad $b =$ \blank{1.5cm} \quad Opens: \blank{2cm}

\end{tcolorbox}

\vspace{0.10in}

%----------------------------------------------------------------------
% Tier A
%----------------------------------------------------------------------
\begin{tcolorbox}[colback=white, colframe=black!40,
    title=\textbf{Tier A --- Approaching Proficiency},
    fonttitle=\bfseries, arc=2mm, left=3mm, right=3mm, top=2mm, bottom=2mm]

\textbf{Directions:} Write the standard form equation for each card below.
Show all work.

\vspace{6pt}
\textbf{Card A:} ``An equation that opens to the right with vertex at
$(-1, 2)$ and passes through $(3, 4)$.'' (Hint: substitute the point to find $a$.)

\writelines{3}

\textbf{Card B:} ``An ellipse centered at the origin with horizontal radius 4
and vertical radius 3.''

\writelines{2}

\textbf{Card C:} ``A hyperbola centered at $(1, 0)$ with $a = 3$ and $b = 2$,
opening left and right.''

\writelines{2}

\textbf{Card D:} ``A circle centered at $(-2, 5)$ with radius 3.''

\writelines{2}

\end{tcolorbox}

\vspace{0.10in}

%----------------------------------------------------------------------
% Tier E
%----------------------------------------------------------------------
\begin{tcolorbox}[colback=white, colframe=black!40,
    title=\textbf{Tier E --- Extension},
    fonttitle=\bfseries, arc=2mm, left=3mm, right=3mm, top=2mm, bottom=2mm]

\textbf{Part 1.} For Card B (Tier A), find all $x$-intercepts and
$y$-intercepts of the ellipse $\dfrac{x^2}{16} + \dfrac{y^2}{9} = 1$.
Show all work.

\writelines{3}

\textbf{Part 2.} For Card C (Tier A), write both asymptote equations for the
hyperbola $\dfrac{(x-1)^2}{9} - \dfrac{y^2}{4} = 1$.

\writelines{2}

\textbf{Part 3.} Create your own Card E. Choose a conic type, state its key
features in words (as if writing a card description), then write its equation.

\vspace{4pt}
My Card E (description): \writelines{2}

My Card E (equation): \writelines{2}

\end{tcolorbox}

\end{document}
